Hace falta conocer la temperatura y si el calefactor está encendido o apagado. Con \(t\) en segundos, sean
\[
t_1=60\log(3/2),\qquad
t_2=60\log2,\qquad P=t_1+t_2=60\log3.
\]
Durante el primer ciclo,
\[
T(t)=
\begin{cases}
60-30e^{-t/60},&0\le t<t_1,\\
20+20e^{-(t-t_1)/60},&t_1\le t<P,
\end{cases}
\]
en grados Celsius; en \(t=P\) vuelve a \(30\,{}^\circ\mathrm C\) y se enciende de nuevo. La evolución se repite con periodo \(P\). La fracción encendida es \(\log(3/2)/\log3\), y
\[
\overline T=20+40\frac{\log(3/2)}{\log3}\quad{}^\circ\mathrm C.
\]
En \(T=35\,{}^\circ\mathrm C\), la derivada temporal es \(5/12\,\mathrm{K/s}\) con el calefactor encendido y \(-1/4\,\mathrm{K/s}\) con él apagado.
